% source: manti.01.ma, 1ère SM, série suites (Drive 0B6pfSJcEVng0ZURtZUk0c3pobk0). Ex. 5 = série française suites_s1 Ex. 4 (non repris)
\ex{Suites}{ZURt}{Ex. 1}{6}
Soit $(U_n)_{n\in\mathbb{N}^*}$ la suite définie par $U_n=\dfrac{n^2+1}{2n^2+n+2}$.
Montrer que $(U_n)_{n\in\mathbb{N}^*}$ est majorée par $\dfrac12$ et minorée par $\dfrac25$.
\fin
% Ex. 2 : même suite que bjRL Ex. 1 (non repris)
\ex{Suites}{ZURt}{Ex. 3}{5}
Soit $(U_n)_{n\in\mathbb{N}}$ la suite définie par $U_0=2$ et $U_{n+1}=\dfrac12U_n+2$.
\begin{enumerate}
\item Calculer $U_1$ et $U_2$.
\item Montrer que $(U_n)_{n\in\mathbb{N}}$ est majorée par $4$.
\item Étudier la monotonie de la suite $(U_n)_{n\in\mathbb{N}}$.
\end{enumerate}
\fin
\ex{Suites}{ZURt}{Ex. 4}{5}
Soit $(U_n)_{n\in\mathbb{N}}$ la suite définie par $U_0=\dfrac32$ et $U_{n+1}=\dfrac{2U_n}{U_n+1}$.
\begin{enumerate}
\item Calculer $U_1$ et $U_2$.
\item Montrer que $(\forall n\in\mathbb{N})\ U_n>1$.
\item Étudier la monotonie de la suite $(U_n)_{n\in\mathbb{N}}$.
\end{enumerate}
\fin
\ex{Suites}{ZURt}{Ex. 6}{9}
Soit $(U_n)_{n\in\mathbb{N}}$ la suite définie par $U_0=-2$ et $U_{n+1}=\dfrac{2U_n-3}{U_n+6}$.
\begin{enumerate}
\item Montrer par récurrence que $(\forall n\in\mathbb{N})\ -3<U_n<-1$.
\item Étudier la monotonie de la suite $(U_n)_{n\in\mathbb{N}}$.
\item On pose $V_n=\dfrac{U_n+1}{U_n+3}$.
\begin{enumerate}
\item Montrer que $(V_n)_{n\in\mathbb{N}}$ est une suite géométrique et déterminer $V_n$ en fonction de $n$.
\item En déduire que $U_n=-\dfrac{5^n+3^{n+1}}{5^n+3^n}$.
\item Calculer en fonction de $n$ la somme $S_n=V_0+V_1+\dots+V_{n-1}$.
\end{enumerate}
\end{enumerate}
\fin
\ex{Suites}{ZURt}{Ex. 7}{10}
Soit $(U_n)_{n\in\mathbb{N}}$ la suite définie par $U_0=12$, $U_1=\dfrac{11}{2}$ et $6U_{n+2}=5U_{n+1}-U_n$.
\begin{enumerate}
\item Calculer $U_2$ et $U_3$.
\item On pose $W_n=3U_{n+1}-U_n$. Montrer que $(w_n)_{n\geq0}$ est une suite géométrique en précisant sa raison, puis déterminer le terme général $W_n$ en fonction de $n$.
\item Montrer que $(\forall n\in\mathbb{N})\ U_{n+1}=\dfrac13U_n+\dfrac{3}{2^{n+1}}$.
\item On pose $V_n=U_n-\dfrac{9}{2^n}$. Montrer que $(V_n)_{n\in\mathbb{N}}$ est une suite géométrique, déterminer $V_n$ en fonction de $n$ et en déduire $U_n$ en fonction de $n$.
\end{enumerate}
\fin
\ex{Suites}{ZURt}{Ex. 8}{7}
Soit $(U_n)_{n\in\mathbb{N}}$ la suite définie par $U_0=1$ et $U_{n+1}=\dfrac14\left(2U_n+n+2\right)$. On pose $V_n=2U_n-n$.
\begin{enumerate}
\item Calculer $U_1$, $V_0$ et $V_1$.
\item Montrer que $(V_n)_{n\in\mathbb{N}}$ est une suite géométrique en précisant sa raison $q$.
\item Calculer le terme général $V_n$ en fonction de $n$, puis en déduire $U_n$ en fonction de $n$.
\end{enumerate}
\fin
